\documentclass{article}

\usepackage{booktabs}
\usepackage{tabularx}
\setlength{\parindent}{0pt}
\setlength{\parskip}{0.75em}

\title{Development Plan\\\progname}

\author{\authname}

\date{}

\input{../Comments.text}
\input{../Common.text}

\begin{document}

\maketitle

\begin{table}[hp]
\caption{Revision History} \label{TblRevisionHistory}
\begin{tabularx}{\textwidth}{llX}
\toprule
\textbf{Date} & \textbf{Developer(s)} & \textbf{Change}\\
\midrule
Sep 27 2026 & Vikram Chandar & Add Expected Technology, Coding Standards and Use of AI sections\\
Sep 27 2026 & Sankalpa Chhetri Dhakal & Add IP to Protect section and Project Decomposition and Scheduling sections\\
Date3 & Name(s) & Description of changes\\
... & ... & ...\\
\bottomrule
\end{tabularx}
\end{table}

\newpage{}

\wss{Put your introductory blurb here.  Often the blurb is a brief roadmap of
what is contained in the report.}

\wss{Additional information on the development plan can be found in the
\href{https://gitlab.cas.mcmaster.ca/courses/capstone/-/blob/main/Lectures/L02b_POCAndDevPlan/POCAndDevPlan.pdf?ref_type=heads}
{lecture slides}.}

\section{Confidential Information?}

\wss{State whether your project has confidential information from industry, or
not.  If there is confidential information, point to the agreement you have in
place.}

\wss{For most teams this section will just state that there is no confidential
information to protect.}
\section{IP to Protect}

There is currently no intellectual property agreement associated with this
project. The project is being developed by the team as part of the capstone
course and is not subject to any additional IP protection requirements from an
external organization.

\section{Copyright License}

\wss{What copyright license is your team adopting.  Point to the license in your
repo.}

\section{Team Meeting Plan}

\wss{This section is for meetings between the team members, and meetings between 
the team and their supervisor/advisor (as appropriate).}

\wss{How often will the team meet? where? when?}

\wss{If the meeting is a physical location (not virtual), out of an abundance of
caution for safety reasons you shouldn't put the location online}

\wss{To help schedule meetings, you can use the lecture/tutorial time when we do 
not have anything else scheduled.}

\wss{How often will you meet with your supervisor/advisor?  when?  where?}

\wss{Will meetings be virtual?  At least some meetings should likely be
in-person.}

\wss{How will the meetings be structured?  There should be a chair for all meetings.  
There should be an agenda for all meetings.}

\section{Team Communication Plan}

\wss{Issues on GitHub should be part of your communication plan.}

\section{Team Member Roles}

\wss{You should identify the types of roles you anticipate, like notetaker,
leader, meeting chair, reviewer.  Assigning specific people to those roles is
not necessary at this stage.  In a student team the role of the individuals will
likely change throughout the year.}

\section{Workflow Plan}

\begin{itemize}
	\item How will you be using git, including branches, pull request, etc.?
	\item How will you be managing issues, including template issues, issue
	classification, etc.?
  \item Use of CI/CD
\end{itemize}

\section{Project Decomposition and Scheduling}

The team will use GitHub Projects to organize and track the development of the
project. A Kanban board has been created with the following stages:
\textit{Backlog}, \textit{Ready}, \textit{In Progress}, \textit{In Review}, and
\textit{Done}. Tasks will move through these stages as work is planned,
implemented, reviewed, and completed.

\vspace{0.75em}

The project is decomposed into smaller GitHub Issues representing development
tasks, research activities, meetings, and course deliverables. Labels are used
to identify the type of work, such as research, development plan tasks, problem
statement tasks, and meetings. This allows the team to assign work to individual
members and monitor the progress of each task.

The GitHub Project can be found at:

\begin{center}
\url{https://github.com/orgs/SoftwareMandems/projects/1}
\end{center}

The overall project schedule will follow the major capstone milestones and
deadlines provided in the course outline. In addition to course deliverables,
technical work will be scheduled in stages. The early stages will focus on setting up the Unity and VR environment and developing a basic FPV drone flight model.
Once the core simulation is working, development will move toward controller
support, training environments, scoring and feedback systems, multiplayer
functionality, and the planned AI-based flight analysis features.

Tasks that are not yet ready to be worked on will remain in the backlog, while
tasks that are prepared for development will move into the ready state. Work
currently being completed will be tracked as in progress, and completed work
will be reviewed before being moved to done. This workflow will allow the team
to continuously track progress and adjust the schedule as priorities change.


\section{Proof of Concept Demonstration Plan}

What is the main risk, or risks, for the success of your project?  What will you
demonstrate during your proof of concept demonstration to convince yourself that
you will be able to overcome this risk?

\section{Expected Technology}


\noindent The project will use an open source game engine like Unity as the primary game engine and C\# as the primary programming language. C\# will be used to implement the drone simulation, flight controls, physics, user interfaces, VR interactions, multiplayer functionality, and AI-based flight analysis. Unity and C\# were selected because they provide established tools for VR development, physics simulation, and cross-platform deployment. However, Unity's abstractions may limit access to lower-level hardware functionality and may introduce performance overhead for computationally intensive components. C may therefore be investigated for lower-level components such as drivers if direct hardware access or performance requirements make it necessary.

\vspace{0.75em}

\noindent Unity's built-in physics, input, and rendering systems will be used where appropriate to support the drone simulation. These systems reduce the need to implement common functionality from scratch, but their general-purpose implementations may not provide the level of control or performance required for realistic FPV flight. Unity's VR and XR packages such as OpenXR Plugin and XR Interaction Toolkit will be investigated for integration with the project's VR hardware, with compatibility and input latency evaluated during development. Additional Unity packages and libraries such as Unity Multiplayer SDK, Unity Transport and Unity Netcode may be used for multiplayer networking, data collection, visualization, and other functionality required by the project. These packages will be evaluated based on their compatibility, performance, documentation, and ability to meet the project's requirements, with custom implementations considered where existing functionality is insufficient.

\vspace{0.75em}

\noindent For the AI-based flight analysis component, pre trained models are expected to be leveraged rather than developing a model from scratch. Open Source models like Qwen2-VL may be investigated to provide feedback on areas such as control efficiency and flight performance. However, general-purpose models may not reliably interpret the structured controller inputs and flight telemetry required by the project and may produce inconsistent feedback. Existing models and libraries will therefore be tested against representative flight data to determine whether they satisfy the project's accuracy, performance, and hardware requirements before considering fine-tuning, adapting an existing model, or developing a custom model.

\vspace{0.75em}

\noindent The project can use a standard linter such as Microsoft C\# linter and formatter to maintain consistent coding conventions and identify potential issues during development. The linter will provide automated checks for common coding problems, although static analysis cannot detect all runtime or simulation-specific issues. Unit testing will therefore be performed using an appropriate C\# testing framework, such as NUnit or Unity's Test Framework. Code coverage tools will also be investigated to measure the extent to which automated tests exercise the implemented functionality and identify areas requiring additional tests.
\vspace{0.75em}

\noindent Continuous Integration (CI) will be implemented using GitHub Actions where practical. Automated builds, tests, and code-quality checks will be run when changes are pushed to the repository or submitted through pull requests. GitHub Actions was selected because it integrates with the team's existing GitHub workflow and can automate repetitive verification tasks. However, Unity projects can require significant build dependencies and execution time, so the CI workflow will be limited to checks that can be reliably automated within the available resources. This will help identify integration and testing issues before changes are merged into the main branch.

\vspace{0.75em}
\noindent Git, GitHub, and GitHub Projects will be used for version control, collaboration, issue tracking, and project management. GitHub Issues will be used to track development tasks and defects, while GitHub Projects will be used to organize and monitor the project's overall progress. These tools provide an integrated workflow for coordinating development, while working within the github ecosystem.

\vspace{0.75em}
\noindent The team will use an open source IDE like Visual Studio Code as the primary IDE for development, debugging, and integration with Unity. Visual Studio Code provides a lightweight development environment and supports the project's required languages and tools, although its Unity and C\# functionality may not provide the same level of integrated tooling as a full-featured IDE but it gives familiarity and flexibility if changes are to be made in the future. Code coverage can be measured using a separate Unity's Code Coverage package to determine the percentage of C\# code executed by the project's automated tests. 

\vspace{0.75em}

\noindent The team will use a free performance profiler such as the Unity Profiler to measure CPU, GPU, memory, physics, and rendering performance throughout development. The Unity Memory Profiler and Frame Debugger will also be investigated to identify memory usage issues and rendering bottlenecks, particularly for maintaining the performance requirements of the VR simulation. These tools provide detailed information about Unity's runtime behaviour, but profiling results may not fully represent performance on the final VR hardware. Performance testing will therefore also be performed on the target hardware to validate that optimizations identified through profiling translate into acceptable real-world performance.

\vspace{0.75em}

\noindent Although existing Unity systems, libraries, and frameworks will be used where appropriate, major project-specific components such as the drone flight model, controller behaviour, flight evaluation system, and integration between the simulation, VR controls, multiplayer functionality, and AI analysis will be implemented through custom code. Existing tools will be preferred when they provide reliable and maintainable functionality, while custom implementations will be considered when their limitations prevent them from meeting the project's requirements. 




\section{Use of AI and LLMs}

\noindent The team will use Claude Code and ChatGPT as development and documentation assistance tools throughout the project. These tools may be used to assist with generating code, explaining unfamiliar concepts, debugging, developing test cases, reviewing code, and drafting or improving project documentation.

\vspace{0.75em}

\noindent LLM-generated code and documentation will not be treated as the standard. All generated code will be reviewed and tested by team members before being incorporated into the project. The team will verify functionality through unit tests, integration testing, manual testing, and performance testing where appropriate. Documentation generated with the assistance of LLMs will also be reviewed by team members against project requirements and rubrics to ensure that it is accurate.

\section{Coding Standard}

The project will use the \textbf{Microsoft C\# Coding Conventions} as the coding standard for all C\# source code. This will provide consistent conventions for naming, formatting, organization, and documentation throughout the project.

\vspace{1em}

\noindent PascalCase will be used for classes, methods, and public properties, and camelCase for local variables and parameters. Methods and classes will have descriptive names, and code will be organized into modular classes with each class having a clear responsibility.

\vspace{1em}

\noindent The project will use consistent indentation and brace formatting, avoid unnecessary code duplication, and include comments where they help explain non-obvious implementation details. Code will be formatted consistently using the development environment's formatting tools. Pull requests will be reviewed by another team member before being merged into the main branch to ensure the coding standard is followed.

\newpage{}

\section{Appendix --- Generative AI Use Disclosure}

\subsection{AI Tools Used}
\begin{itemize}
    \item ChatGPT 5.6 Sol, OpenAI
\end{itemize}

\subsection{How and Where Used}

ChatGPT 5.6 Sol was used only to review the document for grammar, spelling and other
minor phrasing issues. It was not used to generate any technical content, project 
ideas, requirements or diagrams.

Grammar checking was performed throughout the document where necessary to improve
readability and correct grammar errors. None of the 13 sections of this document 
were generated or directly altered by the AI tool.

\subsection{Verification of AI-Generated Content}
All grammar and spelling corrections that were suggested by ChatGPT 5.6 Sol, were
reviewed and added into this document manually. The team ensured that grammatical 
corrections did not alter the intended technical meaning or introduce new factual
or technical claims.

The team remains responsible for the accuracy, completeness, consistency, and 
appropriateness of the submitted document, as no content was directly generated 
by an AI system.

\newpage{}

\section*{Appendix --- Reflection}

\wss{Not required for CAS 741}

\input{../Reflection.text}

\begin{enumerate}
    \item Why is it important to create a development plan prior to starting the
    project?
    \item In your opinion, what are the advantages and disadvantages of using
    CI/CD?
    \item What disagreements did your group have in this deliverable, if any,
    and how did you resolve them?
\end{enumerate}

\newpage{}

\section*{Appendix --- Team Charter}

\wss{borrows from
\href{https://engineering.up.edu/industry_partnerships/files/team-charter.pdf}
{University of Portland Team Charter}}

\subsection*{External Goals}

\wss{What are your team's external goals for this project? These are not the
goals related to the functionality or quality fo the project.  These are the
goals on what the team wishes to achieve with the project.  Potential goals are
to win a prize at the Capstone EXPO, or to have something to talk about in
interviews, or to get an A+, etc.}

\subsection*{Attendance}

\subsubsection*{Expectations}

\wss{What are your team's expectations regarding meeting attendance (being on
time, leaving early, missing meetings, etc.)?}

\subsubsection*{Acceptable Excuse}

\wss{What constitutes an acceptable excuse for missing a meeting or a deadline?
What types of excuses will not be considered acceptable?}

\subsubsection*{In Case of Emergency}

\wss{What process will team members follow if they have an emergency and cannot
attend a team meeting or complete their individual work promised for a team
deliverable?}

\subsection*{Accountability and Teamwork}

\subsubsection*{Quality} 

\wss{What are your team's expectations regarding the quality
of team members' preparation for team meetings and the quality of the
deliverables that members bring to the team?}

\subsubsection*{Attitude}

\wss{What are your team's expectations regarding team members' ideas,
interactions with the team, cooperation, attitudes, and anything else regarding
team member contributions?  Do you want to introduce a code of conduct?  Do you
want a conflict resolution plan?  Can adopt existing codes of conduct.}

\subsubsection*{Stay on Track}

\wss{What methods will be used to keep the team on track? How will your team
ensure that members contribute as expected to the team and that the team
performs as expected? How will your team reward members who do well and manage
members whose performance is below expectations?  What are the consequences for
someone not contributing their fair share?}

\wss{You may wish to use the project management metrics collected for the TA and
instructor for this.}

\wss{You can set target metrics for attendance, commits, etc.  What are the
consequences if someone doesn't hit their targets?  Do they need to bring the
coffee to the next team meeting?  Does the team need to make an appointment with
their TA, or the instructor?  Are there incentives for reaching targets early?}

\subsubsection*{Team Building}

\wss{How will you build team cohesion (fun time, group rituals, etc.)? }

\subsubsection*{Decision Making} 

\wss{How will you make decisions in your group? Consensus?  Vote? How will you
handle disagreements? }

\end{document}